\caption{Lag-dependent response geometry and sampling limitations. (a) The population continuous-time criterion $R_A$ for four exactly sampled, stable two-dimensional real-spectrum systems with a fixed spectrum and varied eigenvector parameter $\kappa\in\{0.3,1,3,6\}$. The computed $R_A$ is invariant over $\tau=1,2,4,8$ to numerical precision while the dotted horizontal line marks $R_A=1$. These results concern the exact two-dimensional generator on a valid logarithm branch, not arbitrary higher-dimensional compressions. (b) The median largest principal angle $\theta_{\max}(Q_\tau,Q_1)$ between the known-operator M2-selected planes at lag $\tau$ and lag one in the $N=40$ embedded benchmark, for changes in spectrum with fixed geometric construction and changes in geometry with fixed spectrum. (c) Percentage of admissible inferred $\widehat R_A$ values for $M=1$ and $M=25$ trajectories in the $N=40$ benchmark. The denominators are 32 fitted instances at each lag and $M$; inadmissible ratios are counted as such, not assigned a numerical value. The settings are reused across lags and must not be treated as independent samples. (d) Complex-frequency aliasing for a known angular frequency $\omega_0=0.72$ per time unit. The principal matrix logarithm recovers the physical frequency before the Nyquist branch boundary $\tau=\pi/\omega_0\approx4.36$, but yields a smaller apparent frequency at $\tau=8$. The aliasing shown is a sampling/branch limitation rather than fitting error. All plotted values are reproduced from the archived V19 T03 test outputs.}
